\section{Examples and comparison to SHAP and Causal SHAP}\label{sec:shap_examples}

\paragraph{Where we are.} Section~\ref{sec:causal_impact} defined \ourapproach
and verified, on Alice and Bob, that the resulting attribution satisfies all seven
desiderata of that section, where PN/PS/PNS fail at least two. The natural next
question is how \ourapproach compares to the dominant practical alternative for
feature attribution in the ML literature --- SHAP and Causal SHAP --- since these
are the tools practitioners typically reach for. This section runs that comparison
on two examples, recording where each method passes or fails an explicit
desideratum. We first define SHAP and Causal SHAP formally
(Section~\ref{sec:shap_defs}), then re-run OBCB
(Section~\ref{sec:obcb_shap}) --- where the two SHAP variants coincide because the
feature set is exogenous, so the divergence between PCI and SHAP is the only
contrast to draw. We then introduce a continuous mediation example
(Section~\ref{sec:signal}) where plain and Causal SHAP do diverge --- one
desideratum (D-YX) is fixed by the interventional correction, others (D-YM,
D-MX, D-MXY) are not --- and tabulate the desiderata in
Section~\ref{sec:shap_signal}. Section~\ref{sec:shap_summary} traces every
failure to one of three structural gaps in the SHAP framework that PCI's witness
mechanism, suspect-set distribution $\Gamma_s$, and realised-outcome reference
jointly close.

Shapley values have become one of the most widely used tools for explaining the
predictions of complex machine learning models. Grounded in cooperative game theory,
they provide a principled decomposition of a model's prediction into additive
contributions from individual features.
\citet{janzing2020feature} argued that the conditional expectation in plain SHAP is the
conceptually wrong way to model feature removal --- a dropped feature should be
marginalised under a $do$-intervention, not under conditioning.
\citet{heskes2020causal} propose Causal SHAP, which implements this prescription by
replacing the observational marginal used for features outside the active coalition
with an interventional distribution: when a dropped feature is causally downstream of
a coalition member, its distribution under the intervention differs from its marginal,
and Causal SHAP corrects for this using Pearl's do-calculus.\footnote{%
A different counterfactual line is taken by \citet{sharma2022cfshapley}, whose
CF-Shapley values use a structural causal model to attribute the change in a metric
from a single fixed reference $x^{\text{ref}}$ to the observed $x$:
$\phi_j \propto \sum_{\mathcal{S}\subseteq N\setminus\{j\}}
[f(x_\mathcal{S}, x_j^{\text{ref}}) - f(x_\mathcal{S}, x_j)]$.
This is the closest neighbour to \ourapproach{} in spirit --- single contrastive
scenario, $do$-style intervention --- and differs in three ways.
\emph{(i)} CF-Shapley commits to one reference $x^{\text{ref}}$, while \ourapproach{}
integrates over an alternative-value distribution $\Delta(\mathbf{x})$ and need not pick
a single contrast.
\emph{(ii)} CF-Shapley returns a one-axis contrast
$f(\cdot, x_j^{\text{ref}}) - f(\cdot, x_j)$, while \ourapproach{} reports a joint
$(Y^s, Y^n)$ measure and applies a user-chosen $ci$, decomposing necessity and
sufficiency separately.
\emph{(iii)} CF-Shapley has no witness mechanism, so it cannot break the
overdetermination and undercutting symmetries that \ourapproach{}'s witness set
$\mathbf{W}$ resolves on the OBCB and signal walkthroughs of this section.}
Formal definitions of both are given in Section~\ref{sec:shap_defs}.

Despite this repair, both methods share deeper limitations. First, both plain and
Causal SHAP can assign non-zero responsibility to variables that play no causal role
with respect to the target --- attribution can run against the direction of the causal
graph, a failure the interventional correction does not address. Second, both methods
explain a model's predicted output averaged over the population, not the individual
factual outcome of a specific person: when the goal is \emph{causal responsibility
attribution}, this produces further failures that are not always made explicit in the
literature. The first limitation will surface in our signal example as violations of
D-YX, D-YM, and D-MX; the second will surface in OBCB as the D-A-rank overdetermination
effect, and in the signal example as Instance~2's budget collapse when the prediction
equals its baseline.

We trace these failures through two examples, comparing plain SHAP, Causal SHAP, and
PCI against the desiderata of Section~\ref{sec:motivations} and some further desiderata related to the example we 
will employ. We first return to the
OBCB model: it permits exact analytic computation and the desiderata are already
known. In OBCB the feature set contains only exogenous roots, so plain and Causal
SHAP coincide --- the comparison is short. We then turn to a continuous mediation
example where downstream variables are observable model inputs; there the two methods
diverge, and we examine what Causal SHAP fixes and what it does not.

\subsection{SHAP and Causal SHAP}\label{sec:shap_defs}

In this section $\mathcal{S} \subseteq N$ denotes a SHAP coalition (a subset of features held
to their factual values during evaluation), playing the same role here that
$\mathbf{C} \subseteq \mathbf{S}$ played in Section~\ref{sec:causal_impact}. We use the
calligraphic $\mathcal{S}$ to keep the SHAP coalition typographically distinct from PCI's
suspect set $\mathbf{S}$ on the printed page; the SHAP literature's plain $S$ is the same
object.

\begin{definition}[SHAP]
A cooperative game is a pair $(N, v)$,
where $N = \{1, \ldots, n\}$ is the set of players and $v: 2^N \to \mathbb{R}$
is the \emph{characteristic function} mapping each subset
(coalition) $\mathcal{S} \subseteq N$ to the value that coalition can produce. The Shapley value
\citep{shapley1953} is the unique attribution satisfying four axioms:
\begin{enumerate}
\item \textbf{Efficiency}: $\sum_i \phi_i = v(N) - v(\emptyset)$
\item \textbf{Symmetry}: If $v(\mathcal{S} \cup \{i\}) = v(\mathcal{S} \cup \{j\})$ for all $\mathcal{S}$, then
      $\phi_i = \phi_j$
\item \textbf{Dummy}: If $v(\mathcal{S} \cup \{i\}) = v(\mathcal{S})$ for all $\mathcal{S}$, then $\phi_i = 0$
\item \textbf{Linearity}: For two games $v_1, v_2$: $\phi_i(\alpha_1 v_1 + \alpha_2 v_2)
      = \alpha_1 \phi_i(v_1) + \alpha_2 \phi_i(v_2)$
\end{enumerate}
The unique solution is:
\begin{equation}
\phi_i = \sum_{\mathcal{S} \subseteq N \setminus \{i\}}
\frac{|\mathcal{S}|!\,(|N|-|\mathcal{S}|-1)!}{|N|!}
\bigl[v(\mathcal{S} \cup \{i\}) - v(\mathcal{S})\bigr]
\label{eq:shapley}
\end{equation}
The weight $\frac{|\mathcal{S}|!(|N|-|\mathcal{S}|-1)!}{|N|!}$ equals the probability that, in a uniformly
random ordering of all players, exactly the members of $\mathcal{S}$ precede $i$. Equivalently,
$\phi_i$ is the average marginal contribution of $i$ across all orderings.
\citet{lundberg2017unified} apply Shapley values to explain model predictions by taking
players to be input features, and the characteristic function to be the expected model
output given a subset of features:
\begin{equation}
v(\mathcal{S}) = \mathbb{E}_{X_{\bar{\mathcal{S}}}}\bigl[f(X_{\mathcal{S}} = x_{\mathcal{S}},\, X_{\bar{\mathcal{S}}})\bigr]
\end{equation}
\noindent where $\bar{\mathcal{S}} = N \setminus \mathcal{S}$ and missing features are marginalized over
their marginal distribution.
\end{definition}

\medskip\noindent Causal SHAP replaces the characteristic function with an
interventional expectation \citep{heskes2020causal}:

\begin{definition}[Causal SHAP]\label{def:causalshap}
\begin{equation}
v_{\text{causal}}(\mathcal{S}) = \mathbb{E}\bigl[f(X) \mid do(X_{\mathcal{S}} = x_{\mathcal{S}})\bigr]
= \int P(X_{\bar{\mathcal{S}}} \mid do(X_{\mathcal{S}} = x_{\mathcal{S}}))\, f(X_{\bar{\mathcal{S}}}, x_{\mathcal{S}})\, dX_{\bar{\mathcal{S}}}
\label{eq:causal_v}
\end{equation}
By the rules of do-calculus, intervening on $X_{\mathcal{S}}$ only affects the distribution of
descendants of $X_{\mathcal{S}}$. For non-descendants, $P(X_j \mid do(X_{\mathcal{S}} = x_{\mathcal{S}})) = P(X_j)$.
Writing $X_{\bar{\mathcal{S}},\,\mathrm{nd}}$ for the features in $\bar{\mathcal{S}}$ that are not
descendants of any member of $\mathcal{S}$, and $X_{\bar{\mathcal{S}},\,\mathrm{d}}$ for those that are,
equation~\eqref{eq:causal_v} is equivalent to:
\begin{equation}
v_{\text{causal}}(\mathcal{S}) = \mathbb{E}_{X_{\bar{\mathcal{S}},\,\mathrm{nd}} \sim P(\cdot),\;
X_{\bar{\mathcal{S}},\,\mathrm{d}} \sim P(\cdot \mid do(X_{\mathcal{S}}=x_{\mathcal{S}}))}\bigl[f(X_{\bar{\mathcal{S}}}, x_{\mathcal{S}})\bigr]
\end{equation}
The Shapley formula~\eqref{eq:shapley} is otherwise unchanged. The method retains all
four Shapley axioms.
\end{definition}

\medskip\noindent
At the level of subset weighting alone, \ourapproach{} (with $J=1$, $K=|\mathbf{S}|$) and
SHAP are not distinguishable: $\Gamma$'s subset weights match the Shapley kernel. The
distinction lies in the integrand. SHAP averages marginal contributions using the
observational conditional
$\mathbb{E}[f(\mathbf{X}_{\mathcal{S}} = \mathbf{x}_{\mathcal{S}}, \mathbf{X}_{\bar{\mathcal{S}}})]$, marginalising over the
empirical distribution of non-coalition features; a feature merely correlated with a true
cause can therefore receive positive attribution. \ourapproach{} evaluates counterfactual
outcomes
$P(Y \in A \mid \mathbf{u},\, \mathrm{do}(\mathbf{C}=\mathbf{c}',\,\mathbf{T}=\mathbf{t}^\star))$
computed via the structural model, so attribution flows only along genuine causal paths.
The witness mechanism adds a further distinction: by holding intermediate variables fixed,
\ourapproach{} can resolve overdetermination and undercutting scenarios where SHAP assigns
equal or incorrect attribution. The remainder of this section makes both points concrete on
two examples.

\subsection{OBCB: SHAP and Causal SHAP}\label{sec:obcb_shap}

To apply the SHAP and Causal SHAP definitions of Section~\ref{sec:shap_defs}
to OBCB we must specify both a feature set $N$ and a function $f$ on the joint
feature space; the value functions and Shapley sums are derived from these.
The natural choice for $f$ is the conditional expectation of the outcome,
$f(x_N)=\mathbb{E}[\varname{loan}\mid X_N=x_N]$. But the OBCB PSCM has five
variables --- $\varname{gender}$, $\varname{credit}$, $\varname{check}$,
$\varname{check\text{-}failed}$, $\varname{loan\_if\_checked}$ --- so the
choice of $N$ is non-trivial. Two options are natural:
\begin{itemize}
\item \textbf{Option A: two upstream roots.}
$N=\{\varname{gender},\varname{credit}\}$. Both features are exogenous; the
three internal variables ($\varname{check}$, $\varname{check\text{-}failed}$,
$\varname{loan\_if\_checked}$) are integrated out inside $f$. This is the
smallest sufficient feature set and matches the way the model is normally
described: predict approval from gender and credit.
\item \textbf{Option B: two roots plus the mediator.}
$N'=\{\varname{gender},\varname{credit},\varname{check\text{-}failed}\}$.
The mediator becomes an observable model input, the same status it has in the
signal-with-mediation example (\S\ref{sec:signal}). This is the comparison
PCI's witness mechanism naturally invites: PCI accesses
$\varname{check\text{-}failed}$ as a witness, so allowing SHAP to access it as
a feature is the fair foil.
\end{itemize}
We work through both. Option~A is simpler and exposes the qualitative failures
of plain SHAP. Option~B then asks whether admitting the mediator rescues Causal
SHAP.

\medskip\noindent\textbf{Option~A: two upstream roots.}
Here $f(g,c)=P(\varname{loan}=1\mid g,c)$ takes values
$f(\text{F},\text{bad})=0.000$, $f(\text{F},\text{good})=0.180$,
$f(\text{M},\text{bad})=0.045$, $f(\text{M},\text{good})=0.900$, with the three
internal variables integrated out. We apply SHAP to the rejection probability
$g(x)=1-f(x)$ for Alice and Bob from Example~\ref{ex:obcb_stochastic},
attributing responsibility for the denial outcome.\footnote{Because $g=1-f$,
the characteristic functions satisfy $v_g(\mathcal{S})=1-v_f(\mathcal{S})$ and
$\phi^g_i=-\phi^f_i$; we compute $\phi^f$ and flip signs. With $|N|=2$ and
uniform marginals $P(\varname{gender})=P(\varname{credit})=\tfrac{1}{2}$,
equation~\eqref{eq:shapley} reduces to
$\phi_i = \tfrac{1}{2}[v(\{i\})-v(\emptyset)] + \tfrac{1}{2}[v(N)-v(N\!\setminus\!\{i\})]$.}
By the SHAP definition (and equivalently by Definition~\ref{def:causalshap}
once we note that no dropped feature is a descendant of any coalition member),
the value function on the 2-feature game is
\[
v(\mathcal{S}) \;=\; \mathbb{E}\!\left[f(X)\,\big|\,X_{\mathcal{S}} = x_{\mathcal{S}}^\star\right]
\;=\; \sum_{x_{N\setminus S}\in\{0,1\}^{|N\setminus S|}} P(X_{N\setminus S}=x_{N\setminus S})\,f(x_{\mathcal{S}}^\star,\,x_{N\setminus S}).
\]
Under independent uniform marginals
$P(\varname{gender})=P(\varname{credit})=\tfrac{1}{2}$ this specialises to
\begin{align*}
v(\emptyset)                              &= \tfrac{1}{4}\sum_{g,c\in\{0,1\}} f(g,c),
&&[\text{both dropped: weight } \tfrac{1}{2}\cdot\tfrac{1}{2}] \\
v(\{\varname{gender}\})                   &= \tfrac{1}{2}\!\left[f(g^\star,0)+f(g^\star,1)\right],
&&[\text{only credit dropped: weight } P(c) = \tfrac{1}{2}] \\
v(\{\varname{credit}\})                   &= \tfrac{1}{2}\!\left[f(0,c^\star)+f(1,c^\star)\right],
&&[\text{only gender dropped: weight } P(g) = \tfrac{1}{2}] \\
v(\{\varname{gender},\varname{credit}\})  &= f(g^\star,c^\star).
&&[\text{nothing dropped: no average}]
\end{align*}
Plugging in the four values of $f$ listed at the start of Option~A yields
Table~\ref{tab:vS_optionA}; pushing those through the closed-form Shapley above
gives the magnitudes in Table~\ref{tab:shap_obcb}, which compares against PNS
and PCI.

\begin{table}[h]
\centering
\begin{tabular}{lcc}
\toprule
$\mathcal{S}$ & $v(\mathcal{S})$ for Alice & $v(\mathcal{S})$ for Bob \\
\midrule
$\emptyset$                              & $0.281$ & $0.281$ \\
$\{\varname{gender}\}$                   & $0.090$ & $0.473$ \\
$\{\varname{credit}\}$                   & $0.023$ & $0.023$ \\
$\{\varname{gender},\varname{credit}\}$  & $0.000$ & $0.045$ \\
\bottomrule
\end{tabular}
\caption{Coalition values $v(\mathcal{S})=\mathbb{E}[f(X)\mid X_{\mathcal{S}}=x_{\mathcal{S}}^\star]$ for the
2-feature game ($N=\{\varname{gender},\varname{credit}\}$).
$v(\emptyset)=\mathbb{E}[f(X)]=0.281$ is the population baseline; $v(N)$ is
the factual prediction $f(x^\star)$. Plain SHAP and Causal SHAP coincide on
this game (neither feature is downstream of the other), so a single column
suffices per individual.}
\label{tab:vS_optionA}
\end{table}

\begin{table}[h]
\centering
\begin{tabular}{llcccc}
\toprule
Person & Feature & PNS & SHAP & PCI (no witnesses) & PCI (with witnesses) \\
\midrule
Alice & \varname{gender} & 0.045 & \phantom{$-$}0.107 & 0.210 & \textbf{0.263} \\
Alice & \varname{credit} & 0.180 & \phantom{$-$}0.174 & 0.240 & \textbf{0.199} \\
Bob   & \varname{gender} & 0.000 & $-$0.107           & 0.038 & \textbf{0.019} \\
Bob   & \varname{credit} & 0.860 & \phantom{$-$}0.343 & 0.228 & \textbf{0.119} \\
\bottomrule
\end{tabular}
\caption{Responsibility attributions for Alice and Bob. SHAP values are signed
($\phi^g_i$ for $P(\varname{loan}=0\mid x)$); desiderata compare magnitudes
$|\cdot|$. PNS and PCI values from Table~\ref{tab:pci_decomp}.}
\label{tab:shap_obcb}
\end{table}

SHAP fails two desiderata. The primary failure is \textbf{D-A-rank}: Alice's 
credit outranks her gender ($0.174 > 0.107$). This is an overdetermination
effect: \varname{credit}=bad correlates with the \varname{check\text{-}failed}
mechanism that blocked Alice's credit evaluation, so SHAP's observational
marginalization inflates credit attribution above gender's. PCI includes the factual
\varname{check\text{-}failed} value directly, isolating gender's upstream causal
role and recovering the correct ranking. PNS fails \textbf{D-A-rank} for the same reason;
it also fails \textbf{D-B1} by returning exactly zero for Bob's gender (see
Table~\ref{tab:desiderata}).

The second SHAP failure is \textbf{D-comp}: $|\phi^g_{\varname{gender}}|=0.107$
for both Alice and Bob. But gender played a strictly larger causal role in Alice's
rejection than in Bob's. For Alice, \varname{gender}=female was sufficient on its
own to produce rejection --- it blocked the credit evaluation entirely, so her
credit was never a factor. For Bob, \varname{gender}=male was a facilitating
condition: it triggered the credit check, but \varname{credit}=bad was the actual
driver of the denial. SHAP cannot distinguish these because it measures how far
each feature value deviates from the population prediction mean --- female and
male are equidistant deviations in opposite directions, so they receive equal
magnitude regardless of their individual causal weight. This is an instance of a
deeper limitation: SHAP explains predicted outputs, not individual factual
outcomes, which we examine further in Section~\ref{sec:signal}.

Causal SHAP coincides exactly with plain SHAP under Option~A: both features
are exogenous roots, so
$P(\varname{credit}\mid do(\varname{gender}=g))=P(\varname{credit})$ and vice
versa, no dropped feature is a descendant of any coalition member, and
$v_{\mathrm{causal}}(\mathcal{S})=v_{\mathrm{plain}}(\mathcal{S})$ for every coalition. Option~A
therefore reduces to a single SHAP analysis, and the failures
above belong to both methods.

\medskip\noindent\textbf{Option~B: two roots plus the mediator.}
Throughout this subsection we abbreviate $\varname{gender}$, $\varname{credit}$, and
$\varname{check\text{-}failed}$ as $g$, $c$, and $cf$ in mathematical expressions and table
cells; running prose retains the full names.
Now we admit \varname{check\text{-}failed} as a third model input and apply
both methods to the 3-feature game $N'=\{\varname{gender},\varname{credit},\varname{check\text{-}failed}\}$.
This requires defining $f$ on $\{0,1\}^3$ rather than on $\{0,1\}^2$. The
natural observational extension is again the conditional expectation
$\tilde f(g,c,cf)=\mathbb{E}[\varname{loan}\mid \varname{gender}=g,\,\varname{credit}=c,\,\varname{check\text{-}failed}=cf]$.
This is well defined for three of the four $(c,cf)$ combinations, but not the
fourth: a failed check requires that a check happened \emph{and} the credit
was bad, so $\varname{check\text{-}failed}=1$ together with
$\varname{credit}=\text{good}$ is impossible under the PSCM and has no
observational data to condition on. We leave $\tilde f$ undefined on this cell.

Both methods need a value function $v(\mathcal{S})$: the expected value of $\tilde f$
when the features in $\mathcal{S}$ are held at their factual values $x_{\mathcal{S}}^\star$ and the
remaining features are averaged out. The methods differ in how that average
is taken. Plain SHAP averages over each dropped feature using its own
observational marginal $P(X_i)$, with the dropped features treated as
independent of one another. Causal SHAP averages over the joint distribution
the PSCM produces under the intervention $do(X_{\mathcal{S}} = x_{\mathcal{S}}^\star)$, so the dropped
features keep whatever joint dependencies the PSCM imposes on them. Whether
either average ever lands on the impossible cell $(c{=}1, cf{=}1)$ --- and so
whether $v(\mathcal{S})$ is finite --- depends on both the method and on which features
are held fixed at which factual values. Table~\ref{tab:vS_optionB} works out
all eight coalition values for both individuals.

\begin{table}[h]
\centering
\small
\begin{tabular}{l cc cc}
\toprule
& \multicolumn{2}{c}{Alice $(g,c,cf){=}(0,0,0)$} & \multicolumn{2}{c}{Bob $(g,c,cf){=}(1,0,1)$} \\
\cmidrule(lr){2-3} \cmidrule(lr){4-5}
$\mathcal{S}$ & $v_{\mathrm{plain}}$ & $v_{\mathrm{causal}}$ & $v_{\mathrm{plain}}$ & $v_{\mathrm{causal}}$ \\
\midrule
$\emptyset$    & \textbf{NaN} & $0.281$ & \textbf{NaN} & $0.281$ \\
$\{g\}$        & \textbf{NaN} & $0.090$ & \textbf{NaN} & $0.473$ \\
$\{c\}$        & $0.007$      & $0.023$ & $0.007$      & $0.023$ \\
$\{cf\}$       & $0.270$      & $0.270$ & \textbf{NaN} & \textbf{NaN} \\
$\{g,c\}$      & $0.000$      & $0.000$ & $0.014$      & $0.045$ \\
$\{g,cf\}$     & $0.090$      & $0.090$ & \textbf{NaN} & \textbf{NaN} \\
$\{c,cf\}$     & $0.000$      & $0.000$ & $0.025$      & $0.025$ \\
$N'$           & $0.000$      & $0.000$ & $0.050$      & $0.050$ \\
\bottomrule
\end{tabular}
\caption{Coalition values for the 3-feature game
$N' = \{\varname{gender},\varname{credit},\varname{check\text{-}failed}\}$.
Bold \textbf{NaN} entries flag coalitions whose value function evaluates
$\tilde f$ at the impossible cell $(c{=}1,cf{=}1)$ with positive weight.
Plain SHAP and Causal SHAP coincide whenever neither $\varname{credit}$ nor
$\varname{check\text{-}failed}$ is dropped, since the only structural coupling
in the model is between these two; they diverge whenever one is in $\mathcal{S}$ and the
other is not.}
\label{tab:vS_optionB}
\end{table}

\textbf{Plain SHAP.}
Plain SHAP treats $\varname{credit}$ and $\varname{check\text{-}failed}$ as
independent, so whenever both are dropped from the coalition the average runs
over all four $(c, cf)$ combinations, with the impossible $(c{=}1, cf{=}1)$
combination receiving weight
$P(c{=}1)\cdot P(cf{=}1) = 0.5\cdot 0.275 \approx 0.140$ ---
the same nonzero weight as any other corner. For Alice, both $\varname{credit}$ and
$\varname{check\text{-}failed}$ are dropped at $\mathcal{S}\in\{\emptyset,\{\varname{gender}\}\}$,
the first two NaN
entries in Table~\ref{tab:vS_optionB}. For Bob the same two coalitions fail,
and two more break for an additional reason: at
$\mathcal{S}=\{\varname{check\text{-}failed}\}$ and
$\mathcal{S}=\{\varname{gender},\varname{check\text{-}failed}\}$ the value $cf^\star=1$
is held fixed while $\varname{credit}$ is sampled from its marginal, so the $c{=}1$ branch
enters with weight $0.5$ and again hits the impossible cell. Plain SHAP gives
undefined Shapley values for both individuals.\footnote{The underlying issue
is that plain SHAP does not see that $\varname{credit}$ and
$\varname{check\text{-}failed}$ are linked by the structural equation
$cf = \varname{check}\cdot(1 - c)$, and so cheerfully samples a combination
the PSCM rules out.}

\textbf{Causal SHAP for Alice.}
Causal SHAP samples $cf$ following its structural equation
$cf = \varname{check}\cdot(1 - c)$, with $\varname{check}\sim\mathrm{Bern}(p_\text{check}[g])$,
given whatever $(g, c)$ values appear in the integrand. When $c=1$ the
equation forces $cf=0$ regardless of $\varname{check}$, so the
$(c{=}1, cf{=}1)$ combination is never produced. Every Alice entry in the
$v_\mathrm{causal}$ column of Table~\ref{tab:vS_optionB} is therefore finite.
Plugging Alice's column into the Shapley formula gives magnitudes
$|\phi_{\varname{gender}}|=0.097$, $|\phi_{\varname{credit}}|=0.176$,
$|\phi_{\varname{check\text{-}failed}}|=0.007$, so
$|\phi_{\varname{credit}}|>|\phi_{\varname{gender}}|$ and \textbf{D-A-rank}
fails: credit outranks gender, exactly the same pathology plain SHAP exhibited
on the 2-feature game.

\textbf{Causal SHAP for Bob.}
Two coalitions break for Bob: $\mathcal{S}=\{\varname{check\text{-}failed}\}$ and
$\mathcal{S}=\{\varname{gender},\varname{check\text{-}failed}\}$. Both hold
$\varname{check\text{-}failed}$ at Bob's factual value $cf^\star=1$ and average
$\varname{credit}$ over its
marginal. To see why this is fatal, apply
Definition~\ref{def:causalshap} to the first:
\[
v_\mathrm{causal}(\{cf\})
\;=\; \mathbb{E}\!\left[\tilde f(X)\,\big|\,do(cf = 1)\right]
\;=\; \tfrac{1}{4}\sum_{g\in\{0,1\}}\sum_{c\in\{0,1\}}\tilde f(g,c,1).
\]
The uniform $\tfrac{1}{4}$ weights come from $\varname{gender}$ and
$\varname{credit}$ sitting upstream of $\varname{check\text{-}failed}$ in the causal graph:
intervening on $\varname{check\text{-}failed}$ overrides the structural equation that would
normally determine it
from credit and check, but does not alter the marginal distributions of the
upstream features. Two of the four summands evaluate $\tilde f$ at
$(g, c{=}1, cf{=}1)$ --- the impossible cell --- with weight $\tfrac{1}{4}$
each. So $v_\mathrm{causal}(\{cf\})$ is undefined, and so is
$v_\mathrm{causal}(\{g, cf\})$ for the same reason. These two coalitions
appear in the Shapley sum for every feature, so all three of Bob's Causal
SHAP values are NaN (Bob's columns in Table~\ref{tab:vS_optionB}).

\textbf{Why Bob breaks and Alice does not.}
The same two coalitions arise for Alice; the difference is the value of
$cf^\star$ being held fixed. With Alice's $cf^\star = 0$, the $c{=}1$ summand
in the formula above evaluates $\tilde f(g, 1, 0) =
p_\text{check}[g]\cdot\text{loan\_prob}[g, 0]$, which is defined because the
PSCM does allow $(c{=}1, cf{=}0)$: when credit is good, $cf$ is forced to $0$
deterministically, so the conditional expectation has a value. With Bob's
$cf^\star = 1$, the $c{=}1$ summand evaluates $\tilde f(g, 1, 1)$, the
impossible cell --- ``good credit and a failed check'' --- which the PSCM
rules out, since failing the credit check requires bad credit by the
structural equation $cf = \varname{check}\cdot(1 - c)$.

The mechanism producing this asymmetry is the intervention itself. When the
PSCM runs forward, $\varname{credit}$ and $\varname{check\text{-}failed}$ are
linked by $cf = \varname{check}\cdot(1 - c)$, which keeps $c{=}1$ and
$cf{=}1$ from co-occurring. When Causal SHAP intervenes via
$do(cf = cf^\star)$, this link is overridden: $\varname{check\text{-}failed}$ is fixed
externally and
$\varname{credit}$ is sampled freely. The override always produces a
$(c{=}1, cf^\star)$ combination with positive weight; whether $\tilde f$ has
a value for that combination depends on which $cf^\star$ value Causal SHAP is
intervening on.

Promoting the mediator to a feature therefore does not rescue Causal SHAP:
it fails the desideratum on Alice and breaks down outright on Bob.

\medskip\noindent
The OBCB analysis exposed two structural failure modes: SHAP's
overdetermination effect on the 2-feature game, and Causal SHAP's breakdown
once the mediator is admitted as a feature under Option~B. The second of
these depends on a structural feature of the OBCB PSCM --- the equation
$cf = \varname{check}\cdot(1 - c)$ ties $\varname{credit}$ and
$\varname{check\text{-}failed}$ together, so Option~B's natural feature set
contains one combination of input values (``good credit and a failed
check'') that the model is not defined on. To see whether the same families
of failure appear without this complication, we now turn to the
signal-with-mediation chain $X \to M \to Y$, where the mediator is naturally
an observable input, $\tilde f$ is defined everywhere, and no impossible-cell
issue arises. Causal SHAP's failures there are not artifacts of an undefined
$\tilde f$ but properties of the method itself.


\subsection{Signal with mediation}\label{sec:signal}

\begin{example}[Signal with mediation]\label{ex:signal_mediation}
We work with the following \textbf{generative model}: $X$ sends a signal,
$M$ mediates it with some noise, and $Y$ is the noisy final outcome:
\begin{align}
X &\sim \mathcal{N}(0.5,\ 0.25) \label{eq:X}\\
M &= X + \varepsilon_M, \quad \varepsilon_M \sim \mathcal{N}(0,\ 0.1) \label{eq:M}\\
Y &= M + \varepsilon_Y, \quad \varepsilon_Y \sim \mathcal{N}(0,\ 0.1) \label{eq:Y}
\end{align}

\noindent The causal graph is $X \to M \to Y$, a simple mediation chain. All noise terms
are independent of each other and of $X$. Since everything is jointly Gaussian, all
conditional expectations are linear functions of the conditioning variables
--- no approximations are required.

Let $R(X \rightsquigarrow Y)$ denote the causal responsibility of $X$ for $Y$.
Following the convention of Section~\ref{sec:causal_impact}, the desiderata are
stated in terms of absolute magnitudes $|R(\cdot)|$, allowing for directionality.
The \textbf{qualitative causal responsibility desiderata} for this example are:
\begin{align*}
\textbf{D-XY} \quad & |R(X \rightsquigarrow Y)| > 0 \\[4pt]
\textbf{D-MY} \quad & |R(M \rightsquigarrow Y)| > 0 \\[4pt]
\textbf{D-XM} \quad & |R(X \rightsquigarrow M)| > 0 \\[8pt]
\textbf{D-YX} \quad & |R(Y \rightsquigarrow X)| = 0 \\[4pt]
\textbf{D-YM} \quad & |R(Y \rightsquigarrow M)| = 0 \\[4pt]
\textbf{D-MX} \quad & |R(M \rightsquigarrow X)| = 0
\end{align*}

A somewhat more subtle desideratum is:
\begin{align*}
\textbf{D-MXY} \quad & |R(M \rightsquigarrow Y)| > |R(X \rightsquigarrow Y)|
\end{align*}
\noindent This captures the intuition that responsibility should be higher for variables
that are closer to the outcome in the causal graph.

\noindent The key \textbf{distributional facts} we will use are:
\begin{align}
\mathbb{E}[X] &= \mathbb{E}[M] = \mathbb{E}[Y] = 0.5, \\
\bigl(\text{Var}(X),\, \text{Var}(M),\, \text{Var}(Y)\bigr) &= (0.25,\, 0.35,\, 0.45), \\
\bigl(\text{Cov}(X,M),\, \text{Cov}(X,Y),\, \text{Cov}(M,Y)\bigr) &= (0.25,\, 0.25,\, 0.35).
\end{align}
\noindent For each target variable, the optimal predictor under squared loss is the
conditional expectation, computable exactly in this Gaussian model:%
\footnote{%
\textbf{Distributional facts.}
From the structural equations:
$\mathrm{Var}(M)=\mathrm{Var}(X)+\mathrm{Var}(\varepsilon_M)=0.25+0.1=0.35$;
$\mathrm{Var}(Y)=\mathrm{Var}(M)+\mathrm{Var}(\varepsilon_Y)=0.35+0.1=0.45$.
$\mathrm{Cov}(X,M)=\mathrm{Cov}(X,X+\varepsilon_M)=\mathrm{Var}(X)=0.25$;
$\mathrm{Cov}(X,Y)=\mathrm{Cov}(X,M+\varepsilon_Y)=\mathrm{Cov}(X,M)=0.25$;
$\mathrm{Cov}(M,Y)=\mathrm{Cov}(M,M+\varepsilon_Y)=\mathrm{Var}(M)=0.35$.
\textbf{Conditional expectations.}
$\mathbb{E}[Y\mid X,M]=M$ since $Y=M+\varepsilon_Y$ with $\varepsilon_Y\perp(X,M)$.
For $\mathbb{E}[M\mid X,Y]$: the covariance matrix of $(X,Y)$ is
$\Sigma_{XY}=\bigl(\begin{smallmatrix}0.25&0.25\\0.25&0.45\end{smallmatrix}\bigr)$
with $\det(\Sigma_{XY})=0.05$, so
$\Sigma_{XY}^{-1}=20\bigl(\begin{smallmatrix}0.45&-0.25\\-0.25&0.25\end{smallmatrix}\bigr)$;
the regression coefficients are
$[\mathrm{Cov}(M,X),\,\mathrm{Cov}(M,Y)]\,\Sigma_{XY}^{-1}
=[0.25,\,0.35]\cdot 20\bigl(\begin{smallmatrix}0.45&-0.25\\-0.25&0.25\end{smallmatrix}\bigr)
=[0.5,\,0.5]$,
giving $\mathbb{E}[M\mid X,Y]=0.5+0.5(X-0.5)+0.5(Y-0.5)=0.5X+0.5Y$.
For $\mathbb{E}[X\mid M,Y]$: by $d$-separation, $X\perp Y\mid M$ in the chain
$X\!\to\!M\!\to\!Y$, so $\mathbb{E}[X\mid M,Y]=\mathbb{E}[X\mid M]
=0.5+\tfrac{\mathrm{Cov}(X,M)}{\mathrm{Var}(M)}(M-0.5)
=0.5+\tfrac{0.25}{0.35}(M-0.5)=0.5+\tfrac{5}{7}(M-0.5)$.}
\begin{align}
f_Y(X, M) &= \mathbb{E}[Y \mid X, M] = M \quad (\{X, M\} \Rightarrow Y) \\
f_M(X, Y) &= \mathbb{E}[M \mid X, Y] = 0.5X + 0.5Y \quad (\{X, Y\} \Rightarrow M) \\
f_X(M, Y) &= \mathbb{E}[X \mid M, Y] = 0.5 + \frac{5}{7}(M - 0.5) \quad (\{M, Y\} \Rightarrow X)
\end{align}

\end{example}


\subsection{Computations and desiderata analysis}\label{sec:shap_signal}

We analyze two instances. \textbf{Instance 1} sets $X^\star=M^\star=Y^\star=1$: every
variable lies above its mean, the prediction is non-trivial at every target, and SHAP,
Causal SHAP, and PCI all return non-zero values. \textbf{Instance 2} sets
$X^\star=M^\star=Y^\star=0.5$: every variable sits at its mean, the prediction equals
its baseline, and SHAP's attribution budget $f(x^\star)-\mathbb{E}[f(X)]$ collapses
to zero. We work through Instance~1 first, then return to Instance~2 to isolate what
each method is really tracking.

\paragraph{Plain SHAP.}
\textit{Target $Y$, features $\{X, M\}$.} First, compute the coalition values:
\begin{align}
v(\emptyset) &= \mathbb{E}[M] = 0.5 \\
v(\{X\}) &= \mathbb{E}[M \mid X=1] = 1 \\
v(\{M\}) &= \mathbb{E}[M \mid M=1] = 1 \\
v(\{X,M\}) &= f_Y(1,1) = 1
\end{align}

Next, apply equation~\eqref{eq:shapley} with $|N|=2$:
\begin{equation}
\phi_X = \tfrac{1}{2}(1 - 0.5) + \tfrac{1}{2}(1 - 1) = 0.25
\qquad
\phi_M = \tfrac{1}{2}(1 - 0.5) + \tfrac{1}{2}(1 - 1) = 0.25
\end{equation}





\textit{Target $M$, features $\{X,Y\}$.}
$v(\emptyset)=0.5$;\enspace $v(\{X\})=1$;\enspace
$v(\{Y\})=0.5\,\mathbb{E}[X\mid Y{=}1]+0.5=0.889$
(using $\mathbb{E}[X\mid Y{=}1]=0.5+\tfrac{0.25}{0.45}(0.5)=0.778$);\enspace
$v(\{X,Y\})=1$.
\begin{align}
\phi_X &= \tfrac{1}{2}(1-0.5)+\tfrac{1}{2}(1-0.889)=0.306, \\
\phi_Y &= \tfrac{1}{2}(0.889-0.5)+\tfrac{1}{2}(1-1)=0.194.
\end{align}

\textit{Target $X$, features $\{M,Y\}$.}
$v(\emptyset)=0.5$;\enspace $v(\{M\})=0.5+\tfrac{5}{7}(0.5)=0.857$;\enspace
$v(\{Y\})=\mathbb{E}[X\mid Y{=}1]=0.778$ (as above);\enspace
$v(\{M,Y\})=\mathbb{E}[X\mid M{=}1]=0.857$ (since $X\perp Y\mid M$).
\begin{align}
\phi_M &= \tfrac{1}{2}(0.857-0.5)+\tfrac{1}{2}(0.857-0.778)=0.219, \\
\phi_Y &= \tfrac{1}{2}(0.778-0.5)+\tfrac{1}{2}(0.857-0.857)=0.139.
\end{align}

\noindent Plain SHAP assigns $\phi_Y = 0.139 > 0$ here, even though $Y$ does not
cause $X$. This arises because $v(\{Y\}) = 0.778 > 0.5$: observing $Y{=}1$ is
statistically informative about $X{=}1$ through the chain $X \to M \to Y$. SHAP
cannot distinguish ``$Y$ is informative about $X$'' from ``$Y$ causes $X$''
because the value function is grounded in observational conditionals, which are
symmetric under reversal of causal direction.



\paragraph{Causal SHAP.}\label{sec:causal_shap_signal}
We now apply Causal SHAP (Definition~\ref{def:causalshap}) and compare against
plain SHAP.

\textit{Target $Y$, features $\{X,M\}$.}
$v_{\mathrm{causal}}(\emptyset)=0.5$;\enspace
$v_{\mathrm{causal}}(\{X\})=\mathbb{E}[M\mid do(X{=}1)]=1$
(M is downstream of X: $P(M\mid do(X{=}1))=\mathcal{N}(1,0.1)$);\enspace
$v_{\mathrm{causal}}(\{M\})=1$ ($do(M{=}1)$ sets $f_Y=M=1$ directly,
$P(X\mid do(M{=}1))=P(X)$ since X has no causal parents);\enspace
$v_{\mathrm{causal}}(\{X,M\})=1$.
All values match plain SHAP, so $\phi_X=\phi_M=0.25$.

\textit{Target $M$, features $\{X,Y\}$.}
$v_{\mathrm{causal}}(\emptyset)=0.5$;\enspace
$v_{\mathrm{causal}}(\{X\})=1$ ($Y$ is downstream of $X$, same as plain SHAP);\enspace
$v_{\mathrm{causal}}(\{Y\})=\mathbb{E}_{X\sim P(X)}[0.5X+0.5{\cdot}1]=0.75$
($Y$ does not cause $X$, so $P(X\mid do(Y{=}1))=P(X)$;
differs from plain SHAP's $0.889$);\enspace
$v_{\mathrm{causal}}(\{X,Y\})=1$.
\begin{align}
\phi_X &= \tfrac{1}{2}(1-0.5)+\tfrac{1}{2}(1-0.75)=0.375, \\
\phi_Y &= \tfrac{1}{2}(0.75-0.5)+\tfrac{1}{2}(1-1)=0.125.
\end{align}

\textit{Target $X$, features $\{M,Y\}$.}
$v_{\mathrm{causal}}(\emptyset)=0.5$;\enspace
$v_{\mathrm{causal}}(\{M\})=0.857$ ($Y$ is downstream of $M$ but $f_X$
does not depend on $Y$ since $X\perp Y\mid M$, so
$\mathbb{E}_{Y\sim P(Y\mid do(M{=}1))}[f_X(1,Y)]=f_X(1,\cdot)=0.857$;
same as plain SHAP);\enspace
$v_{\mathrm{causal}}(\{Y\})=\mathbb{E}_{M\sim P(M)}[f_X(M,1)]=0.5$
($M$ is not a descendant of $Y$, so $P(M\mid do(Y{=}1))=P(M)$
and $\mathbb{E}[f_X(M,1)]=0.5+\tfrac{5}{7}(\mathbb{E}[M]-0.5)=0.5$;
differs from plain SHAP's $0.778$);\enspace
$v_{\mathrm{causal}}(\{M,Y\})=0.857$.
\begin{align}
\phi_M &= \tfrac{1}{2}(0.857-0.5)+\tfrac{1}{2}(0.857-0.5)=0.357, \\
\phi_Y &= \tfrac{1}{2}(0.5-0.5)+\tfrac{1}{2}(0.857-0.857)=0.
\end{align}



\paragraph{PCI computation.}\label{sec:pci_signal_computation}
We work at the factual instance $X^\star = M^\star = Y^\star = 1$ with the
absolute-difference score (Example~\ref{ex:absolute_score})
$ci(y^s, y^n, y^\star) = |y^n - y^\star| - |y^s - y^\star|$. Both worlds run
the PSCM with fresh exogenous noise drawn from $P_{\mathbf{U}}$: the necessity
world replaces the suspect $A$ with $A' \sim P(A)$ and propagates; the
sufficiency world restores it to $a^\star$ and propagates. For target $B$,
$\mathbf{S}$ is the two input features of $B$ and $\mathbf{W}$ the remaining
input; $\Gamma_s$ is uniform over the three non-empty subsets of $\mathbf{S}$
and $\Gamma_w$ uniform over $2^{\mathbf{W}}$.

\textit{$R(X \rightsquigarrow Y)$.}
$\mathbf{S} = \{X, M\}$, $\mathbf{W} = \{M\}$, $w^\star = m^\star = 1$, and
$X' \sim \mathcal{N}(0.5,\, 0.25)$. Three of the four accepted
$(\mathbf{C}, \mathbf{T})$ pairs intervene on $X$, each with weight
$\tfrac{1}{4}$; the fourth, $(\mathbf{C}{=}\{M\}, \mathbf{T}{=}\emptyset)$,
does not touch $X$ and contributes nothing. The key observation is that the
witness pair $(\mathbf{C}{=}\{X\}, \mathbf{T}{=}\{M\})$ pins $M$ at $m^\star$
in both worlds, which equalizes their distributions and zeros out $\bar c$;
the other two pairs let $M$ propagate freely. Computing
$\mathbb{E}[|Y^n - 1|]$ and $\mathbb{E}[|Y^s - 1|]$ in closed form via the
folded-normal formula gives:
\begin{center}
\begin{tabular}{lccc}
\toprule
$(\mathbf{C}, \mathbf{T})$ & $\mathbb{E}[|Y^n - 1|]$ & $\mathbb{E}[|Y^s - 1|]$ & $\bar c$ \\
\midrule
$(\{X\}, \emptyset)$              & $0.677$ & $0.357$ & $0.320$ \\
$(\{X\}, \{M\})$ \emph{[witness]} & $0.252$ & $0.252$ & $0.000$ \\
$(\{X, M\}, \emptyset)$           & $0.677$ & $0.252$ & $0.425$ \\
\bottomrule
\end{tabular}
\end{center}
Combining at weight $\tfrac{1}{4}$ each,
\[
R(X \rightsquigarrow Y \mid \mathbf{W}{=}\{M\})
= \tfrac{1}{4}(0.320 + 0 + 0.425) \approx 0.186.
\]

\textit{$R(Y \rightsquigarrow X)$.}
$X$ is exogenous: no $do$-intervention on $\{M, Y\}$ alters its structural
equation, so necessity and sufficiency distributions coincide and
$R(Y \rightsquigarrow X \mid \mathbf{W}{=}\{M\}) = 0$.

\textit{Other cells.}
The remaining entries of Table~\ref{tab:signal_desid} (D-MY, D-XM, D-YM, D-MX,
D-MXY, and the $\mathbf{W} = \emptyset$ column) follow the same closed-form
machinery; full per-pair computations are in the companion notebook
\texttt{signal\_mediation\_computations.ipynb}.

\paragraph{Desiderata analysis.}\label{sec:signal_desiderata}
Following the convention of Section~\ref{sec:causal_impact}, we use $R(A \rightsquigarrow B)$
for the method-agnostic responsibility and $\phi_A^B$ exclusively for SHAP/Causal SHAP values;
prose statements about desiderata refer to magnitudes $|R(A \rightsquigarrow B)|$.
Table~\ref{tab:signal_desid} maps the desiderata from Section~\ref{sec:signal}
to the values computed above.

\begin{table}[ht]
\centering
\footnotesize
\setlength{\tabcolsep}{4pt}
\resizebox{\textwidth}{!}{%
\begin{tabular}{llcccc}
\toprule
Desideratum & Condition & Plain SHAP & Causal SHAP
  & PCI ($\mathbf{W}{=}\emptyset$) & PCI ($\mathbf{W}{=}$3rd) \\
\midrule
\textbf{D-XY}  & $|R(X \rightsquigarrow Y)| > 0$
  & $0.250$\;\checkmark & $0.250$\;\checkmark
  & $0.249$\;\checkmark & $0.186$\;\checkmark$^\dagger$ \\
\textbf{D-MY}  & $|R(M \rightsquigarrow Y)| > 0$
  & $0.250$\;\checkmark & $0.250$\;\checkmark
  & $0.283$\;\checkmark & $0.319$\;\checkmark \\
\textbf{D-XM}  & $|R(X \rightsquigarrow M)| > 0$
  & $0.306$\;\checkmark & $0.375$\;\checkmark
  & $0.253$\;\checkmark & $0.284$\;\checkmark \\
\textbf{D-YX}  & $|R(Y \rightsquigarrow X)| = 0$
  & $0.139$\;\texttimes & $0.000$\;\checkmark
  & $0.000$\;\checkmark & $0.000$\;\checkmark \\
\textbf{D-YM}  & $|R(Y \rightsquigarrow M)| = 0$
  & $0.194$\;\texttimes & $0.125$\;\texttimes
  & $0.000$\;\checkmark & $0.000$\;\checkmark \\
\textbf{D-MX}  & $|R(M \rightsquigarrow X)| = 0$
  & $0.219$\;\texttimes & $0.357$\;\texttimes
  & $0.000$\;\checkmark & $0.000$\;\checkmark \\
\textbf{D-MXY} & $|R(M \rightsquigarrow Y)| > |R(X \rightsquigarrow Y)|$
  & $0.25{=}0.25$\;\texttimes & $0.25{=}0.25$\;\texttimes
  & $0.283{>}0.249$\;\checkmark & $0.319{>}0.186$\;\checkmark$^\ddagger$ \\
\bottomrule
\end{tabular}%
}
\caption{Desiderata satisfied (\checkmark) and violated (\texttimes) by plain SHAP,
Causal SHAP, and PCI on the signal-with-mediation example
($X\!\to\!M\!\to\!Y$), instance $X=M=Y=1$. PCI uses
$\bar c = \mathbb{E}[|B^n-B^\star|] - \mathbb{E}[|B^s-B^\star|]$ with $\Gamma_s,\Gamma_w$
uniform over non-empty subsets and rejection on
$\mathbf{C}\cap\mathbf{T}\neq\emptyset$. \textbf{W=3rd} sets the third variable
(neither $A$ nor $B$) as the witness when not in $\mathbf{C}$.
$^\dagger$~0.186 weights three suspect-containing $(C,T)$ configurations at
$1/4$ each (the rejection-sampling normalizer $Z=2/3$ also counts a fourth
valid pair $(\{M\},\emptyset)$ that does not contain $X$): one pins $M$ as
witness (contribution~0, necessity and sufficiency coincide) and two let
$M$ propagate freely (contributions~0.320 and 0.425).
$^\ddagger$~The two values use different witness sets ($\mathbf{W}{=}\{M\}$
for D-XY, $\mathbf{W}{=}\{X\}$ for D-MY); this column reports the amplified
verdict under an asymmetric witness rule. The $\mathbf{W}{=}\emptyset$ column
above is the corresponding single-game ranking and already passes D-MXY.}
\label{tab:signal_desid}
\end{table}

\noindent\textbf{Positive results (D-XY, D-MY, D-XM).}
Both methods correctly assign positive attribution to every direct or upstream cause
of the target. Causal prediction models are built from observational conditionals,
which do preserve correlation from cause to effect, so these desiderata pose no
challenge to either method.

\noindent\textbf{Causal SHAP fixes D-YX, not D-YM.}
For target $X$, the interventional value $v_{\mathrm{causal}}(\{Y\})=0.5$ equals
the baseline because $M$ is not downstream of $Y$: $P(M\mid do(Y{=}1))=P(M)$,
so fixing $Y=1$ provides no interventional leverage over $f_X$, and Causal SHAP
correctly yields $\phi_Y^X=0$.
For target $M$, however, $\phi_Y^M=0.125>0$ persists.
The prediction model $f_M(X,Y)=0.5X+0.5Y$ contains $Y$ as an input because $Y$ is
observationally predictive of $M$; setting $Y=1$ raises the model output even under
the interventional regime.
Causal SHAP replaces the conditioning distribution but cannot remove a feature that
the prediction model itself uses.

\noindent\textbf{D-MX worsens under Causal SHAP.}
Plain SHAP assigns $\phi_M^X=0.219$; Causal SHAP assigns $0.357$.
Fixing D-YX reduces $v_{\mathrm{causal}}(\{Y\})$ from $0.778$ to $0.5$, driving
$\phi_Y^X$ to zero, but reallocates the freed attribution to~$M$.
Since the structural noise that actually drives $X$ is unobserved, all attribution
concentrates on the only visible predictor,~$M$, even though $M$ does not cause~$X$.

\noindent\textbf{D-MXY fails for both.}
Since $f_Y=M$, the singleton values $v(\{X\})=\mathbb{E}[M\mid X{=}1]=1$ and
$v(\{M\})=1$ are identical, giving $\phi_M^Y=\phi_X^Y=0.25$ for both methods.
The Shapley formula cannot distinguish that $X$ reaches $Y$ only through~$M$:
when $X=1$, the expected prediction is already~$1$ regardless of whether $M$ is
observed or not, so both features receive the same marginal contribution.

\noindent\textbf{PCI: structural intervention separates the paths; witnesses
sharpen the gap.}
Table~\ref{tab:signal_desid} reports PCI in two configurations: without
witnesses ($\mathbf{W}=\emptyset$) and with the third variable as witness.

Without witnesses, PCI satisfies all six single-variable desiderata.
D-YX, D-YM, and D-MX are zero exactly: intervening on a variable that does not
structurally affect the target leaves both the necessity and sufficiency
worlds with identical distributions, so $ci=0$ identically.
D-XY, D-MY, and D-XM are positive.
D-MXY also \textbf{passes} at $\mathrm{PCI}(M\to Y)\approx 0.283 > \mathrm{PCI}(X\to Y)\approx 0.249$,
but only by a margin of $0.034$. The asymmetry comes from the sufficiency
world: at $C{=}\{X\}$, $X$'s path to $Y$ accumulates
$\varepsilon_M+\varepsilon_Y$ in $|Y^s-Y^\star|$, while at $C{=}\{M\}$, $M$'s
path accumulates only $\varepsilon_Y$. The necessity expectations match
across paths ($\mathrm{Var}(X)+\mathrm{Var}(\varepsilon_M)+\mathrm{Var}(\varepsilon_Y)
=\mathrm{Var}(M)+\mathrm{Var}(\varepsilon_Y)=0.45$), so the gap is carried
entirely by the sufficiency variance asymmetry $\mathrm{Var}(\varepsilon_M)$.
A small structural margin like this is fragile: under any perturbation of
the noise variances or the contrast function the $\mathbf{W}=\emptyset$
verdict could flip.

Admitting the third variable as a witness amplifies the gap by roughly
fourfold.
$\mathrm{PCI}(X\to Y\mid\mathbf{W}{=}\{M\})\approx 0.186$ weights three
suspect-containing $(C,T)$ configurations at $1/4$ each: in two cases $M$
propagates freely (contributing~$0.320$ and $0.425$) and in one case $M$ is
pinned as witness in both worlds, equalizing them and contributing~$0$.
$\mathrm{PCI}(M\to Y\mid\mathbf{W}{=}\{X\})\approx 0.319$ weights three cases
in which $M$ takes an alternative value; since $X$ does not appear in $Y$'s
structural equation, pinning $X$ has no effect on $Y^n$ or $Y^s$, and all
three cases contribute the same $0.425$.
The gap widens from $0.034$ at $\mathbf{W}=\emptyset$ to
$0.319-0.186=0.133$ at $\mathbf{W}=\{$third$\}$ --- D-MXY is now a
\emph{robust} attribution rather than a marginal one.
The comparison relies on different witness sets for the two cells, so it
remains illustrative rather than a single-game ranking.

\paragraph{Instance 2 (baseline).}\label{sec:signal_baseline}
We now return to \textbf{Instance 2}: every variable sits at its mean,
$X^\star = M^\star = Y^\star = 0.5$, and the model predicts the baseline at every
target: $f_Y(0.5,0.5)=0.5$, $f_M(0.5,0.5)=0.5$, $f_X(0.5,0.5)=0.5$. At this
instance, SHAP's attribution budget $f(x^\star)-\mathbb{E}[f(X)]$ vanishes, while
PCI's machinery is unaffected.

\begin{table}[h]
\centering
\small
\begin{tabular}{lcccccc}
\toprule
& \multicolumn{2}{c}{Plain / Causal SHAP} & \multicolumn{2}{c}{PCI ($\mathbf{W}=\emptyset$)} & \multicolumn{2}{c}{PCI ($\mathbf{W}=\{$third$\}$)} \\
\cmidrule(lr){2-3}\cmidrule(lr){4-5}\cmidrule(lr){6-7}
desideratum & value & status & value & status & value & status \\
\midrule
D-XY & $0.000$ & $\times$ & $0.154$ & $\checkmark$ & $0.115$ & $\checkmark$ \\
D-MY & $0.000$ & $\times$ & $0.189$ & $\checkmark$ & $0.212$ & $\checkmark$ \\
D-XM & $0.000$ & $\times$ & $0.146$ & $\checkmark$ & $0.165$ & $\checkmark$ \\
D-YX & $0.000$ & $\checkmark$ & $0.000$ & $\checkmark$ & $0.000$ & $\checkmark$ \\
D-YM & $0.000$ & $\checkmark$ & $0.000$ & $\checkmark$ & $0.000$ & $\checkmark$ \\
D-MX & $0.000$ & $\checkmark$ & $0.000$ & $\checkmark$ & $0.000$ & $\checkmark$ \\
\midrule
D-MXY & \multicolumn{2}{c}{$0.000=0.000$~$\times$} & \multicolumn{2}{c}{$0.189>0.154$~$\checkmark$} & \multicolumn{2}{c}{$0.212>0.115$~$\checkmark$} \\
\bottomrule
\end{tabular}
\caption{Desiderata at the baseline instance $X^\star=M^\star=Y^\star=0.5$.
SHAP collapses to zero on every cell --- correctly assigning zero to the
non-causal triples by accident, but with no signal to offer on the causal
ones. PCI reports the same structural verdicts as at $X^\star=M^\star=Y^\star=1$.}
\label{tab:signal_desid_baseline}
\end{table}

The pattern in Table~\ref{tab:signal_desid_baseline} reflects the object each
method is decomposing. SHAP attributes the quantity
$f(x^\star)-\mathbb{E}[f(X)]$ across features. At the baseline this gap is
zero on every target, so the budget being distributed is zero, and every cell
collapses to zero --- uniformly, regardless of how the realized outcome was
actually produced. The non-causal triples (D-YX, D-YM, D-MX) are satisfied only
by this collapse: SHAP is not isolating the absence of a causal path; it is
reporting that there is nothing to attribute. The same collapse turns D-XY,
D-MY, and D-XM into violations.

PCI never invokes a population expectation. With $\mathbf{S}$ and $\mathbf{W}$
defined relative to a fixed factual instance, the question is whether varying
the suspect would move the realized outcome $y^\star$ away from itself ---
not whether $f(x^\star)$ deviates from $\mathbb{E}[f(X)]$. The structural
mechanism is unchanged at the baseline, so the three causal pairs continue
to register positive values; the three non-causal pairs continue to register
zero exactly, because intervening on a variable that does not structurally
affect the target leaves both worlds with identical distributions. The D-MXY
verdict is preserved in both witness configurations, with margin
$0.035$ at $\mathbf{W}=\emptyset$ (vs.\ $0.034$ at the
$X^\star=M^\star=Y^\star=1$ instance) and $0.097$ at
$\mathbf{W}=\{$third$\}$ (vs.\ $0.133$). The gap is a property of the
PSCM's noise variances, not of where the factual sits.

The two instances together make the contrast sharp. At Instance~1, SHAP's
attribution budget $f(x^\star)-\mathbb{E}[f(X)]$ is non-zero and SHAP returns
non-trivial numbers; at Instance~2 the budget vanishes and every cell collapses.
PCI returns the same structural verdicts on all seven desiderata across the two
instances. The lesson is not that SHAP gives the wrong answer at the baseline ---
it correctly reports that the factual prediction does not deviate from its mean.
The point is that this is a different question from causal responsibility for
$y^\star$: when the prediction equals its baseline, SHAP has no budget to
distribute even when the PSCM's causal mechanism is perfectly intact. PCI asks
the responsibility question directly, and answers it whether or not the
prediction happens to coincide with its population mean.

\subsection{Summary across both examples}\label{sec:shap_summary}

The failures above are not artifacts of a particular feature set or
parameter choice: each traces to a structural limitation of the SHAP/Causal
SHAP framework. Causal SHAP strictly improves on plain SHAP for one
desideratum~(D-YX), by replacing observational conditioning with
interventional distributions; the remaining failures share two sources ---
prediction models include non-causal features that the conditioning
distribution cannot exclude~(D-YM, D-MX), and Shapley averaging cannot
distinguish direct from indirect causal paths~(D-MXY). The OBCB analysis of
Section~\ref{sec:obcb_shap} exhibits the same family in different surface
form: plain SHAP fails \textbf{D-A-rank} on the 2-feature game, and Causal
SHAP either inherits that failure (Alice) or becomes undefined outright
(Bob) once the mediator is admitted as a feature. Across both examples,
neither SHAP variant reliably satisfies the desiderata that motivated
\Ourapproach in Section~\ref{sec:motivations}. PCI addresses these gaps
through three ingredients the SHAP framework systematically lacks --- the
witness mechanism, the suspect-set distribution $\Gamma_s$, and the use of
the realized outcome rather than the prediction baseline; the path-noise
asymmetry in the sufficiency world produces a structural gap on D-MXY,
which appears already in the witness-free single game and is amplified into a
robust margin under a witness rule.

\paragraph{Where we go next.} The two examples in this section are small enough
to inspect every desideratum individually, but they cover only two causal
patterns. Section~\ref{sub:synthetic} stress-tests \ourapproach on a synthetic
structural model carrying three textbook archetypes --- linear
necessary-and-sufficient, overdetermined-not-necessary, and preempted (gated)
--- plus an irrelevance control, and verifies that \ourapproach separates them
at a Monte Carlo budget small enough to bootstrap honest error bars.
Section~\ref{sec:DCE} then contrasts \ourapproach with gradient-based attribution
on a continuous credit-limit example.


